\section{JsonParse.lean}

\code{JsonParse.lean} defines JSON entry points for the structured
\code{Diff} type. The file does not define the JSON schema itself. The schema
shape is implemented by the \code{FromJson Diff} and \code{ToJson Diff}
instances from \code{Defns.lean}. This file gives those instances short,
named wrappers.

\subsection{Overall Structure}

The file has three layers.

\begin{enumerate}
  \item Imports and namespace setup.
  \item Three executable functions:
    \code{parseDiffJson}, \code{parseDiff}, and \code{diffAsJson}.
  \item Two definitional theorems describing exactly what the wrappers expand to.
\end{enumerate}

\subsection{Imports and Namespace}

\begin{definition}[import VersionControl.Defns]
\code{[import VersionControl.Defns]} imports \code{Diff}, \code{FromJson Diff},
\code{ToJson Diff}, and the \code{VersionControl} namespace declarations needed
by this file.
\end{definition}

\begin{definition}[open Lean]
\code{[open Lean]} makes \code{Json}, \code{Json.parse}, \code{fromJson?}, and
\code{toJson} available without writing the \code{Lean.} prefix.
\end{definition}

\begin{definition}[namespace VersionControl]
\code{[namespace VersionControl]} places the parsing functions and theorems in
the same namespace as \code{Diff}. The full names are
\code{VersionControl.parseDiffJson}, \code{VersionControl.parseDiff}, and
\code{VersionControl.diffAsJson}.
\end{definition}

\subsection{JSON and Except}

\begin{definition}[Json]
\code{[Json]} is Lean's datatype for already-parsed JSON values. A value of
type \code{Json} is not raw text. It is a structured JSON tree.
\end{definition}

\begin{definition}[Json.parse]
\code{[Json.parse]} has type \code{String -> Except String Json}. It reads raw
JSON text. If the text is syntactically invalid JSON, it returns
\code{Except.error message}. If the text is syntactically valid JSON, it returns
\code{Except.ok json}.
\end{definition}

\begin{definition}[Except String alpha]
\code{[Except String alpha]} is a result type with two cases:
\code{Except.ok value} for success and \code{Except.error message} for failure.
In this file, parsing functions use \code{Except String} so syntax errors and
schema-conversion errors can carry an error message.
\end{definition}

\begin{definition}[do notation]
\code{[do]} notation sequences computations that can fail. In
\code{parseDiff}, the line \code{let json <- Json.parse text} means: parse the
raw text; if parsing fails, return that error immediately; if parsing succeeds,
bind the parsed JSON value to \code{json} and continue.
\end{definition}

\subsection{Functions}

\begin{definition}[parseDiffJson]
\code{[parseDiffJson]} has type \code{Json -> Except String Diff}. It converts
an already-parsed JSON value into a \code{Diff}. It is exactly
\code{fromJson? json} specialized to \code{Diff}.
\end{definition}

\begin{definition}[fromJson?]
\code{[fromJson?]} is the generic Lean JSON conversion function. In this file,
\code{fromJson? json} uses the \code{FromJson Diff} instance defined in
\code{Defns.lean}, because the expected result type is \code{Diff}.
\end{definition}

\begin{definition}[parseDiff]
\code{[parseDiff]} has type \code{String -> Except String Diff}. It is the main
entry point for raw JSON text. It first calls \code{Json.parse}; then it calls
\code{parseDiffJson} on the parsed JSON value.
\end{definition}

\begin{definition}[diffAsJson]
\code{[diffAsJson]} has type \code{Diff -> Json}. It serializes a \code{Diff}
to a JSON value by calling \code{toJson diff}.
\end{definition}

\begin{definition}[toJson]
\code{[toJson]} is the generic Lean JSON serialization function. In this file,
\code{toJson diff} uses the \code{ToJson Diff} instance defined in
\code{Defns.lean}, because the input value has type \code{Diff}.
\end{definition}

\subsection{Parsing Pipeline}

The raw-text parsing pipeline is:

\begin{enumerate}
  \item Input: \code{text : String}.
  \item \code{Json.parse text} checks JSON syntax.
  \item If JSON syntax fails, \code{parseDiff text} returns the syntax error.
  \item If JSON syntax succeeds, Lean has \code{json : Json}.
  \item \code{parseDiffJson json} checks whether the JSON has the fields and
    types required by \code{Diff}.
  \item If conversion fails, \code{parseDiff text} returns the conversion error.
  \item If conversion succeeds, \code{parseDiff text} returns
    \code{Except.ok diff}.
\end{enumerate}

This file does not call \code{Diff.checkSchema}. Therefore
\code{parseDiff text = Except.ok diff} means the JSON was converted into a
\code{Diff} value. To check hunk shape and pairwise non-overlap, call
\code{diff.checkSchema} after parsing.

\subsection{Theorems}

\begin{theorem}[parseDiff\_eq\_parseDiffJson]
\code{[parseDiff_eq_parseDiffJson]} states that \code{parseDiff text} is exactly
the two-step computation that first runs \code{Json.parse text} and then runs
\code{parseDiffJson json}.
\end{theorem}

\begin{theorem}[parseDiffJson\_eq\_fromJson]
\code{[parseDiffJson_eq_fromJson]} states that \code{parseDiffJson json} is
definitionally equal to \code{fromJson? json} specialized to
\code{Except String Diff}.
\end{theorem}

\subsection{Downstream Use}

\code{JsonParse.lean} feeds into \code{Semantics.lean}. After JSON has been parsed
into a \code{Diff}, the next operation is applying that \code{Diff} to a
\code{BaseFile}, which is described in \code{Semantics.lean}.

\subsection{Integration Note}

The repository build checks this file through \code{VersionControl.lean}.
